% =====================================================================
\section{Start Here}
\label{sec:start}
% =====================================================================

This guide explains all the mathematics behind our project, from the ground up. The base paper proposes a way to
\emph{turn} a rigid body (a satellite, a drone) to a desired orientation while guaranteeing in advance how the motion
will look: how far it may overshoot, how fast it must settle and how small the final error must be, and all of this
with motors that cannot deliver unlimited torque. We implemented that controller, reproduced the paper's
simulations, and then carried it onto a quadcopter (the Gazebo X3).

\subsection{The big picture in plain words}

Controlling the attitude of a body means answering the same five questions many times per second.

\begin{center}\small
\begin{tabularx}{\linewidth}{@{}>{\raggedright\arraybackslash}p{4.6cm} X >{\raggedright\arraybackslash}p{2.4cm}@{}}
\toprule
\textbf{Question} & \textbf{What answers it} & \textbf{Where}\\
\midrule
\emph{How does the body turn at all?} & Physics: quaternion kinematics and Euler's equation for a spinning body & block 1\\
\emph{What can the motors really do?} & A saturation model: commanded torque is clipped at $u_{\max}$ & block 2\\
\emph{How should the motion look?} & A time-varying ``tube'' (envelope) the states must stay inside & block 3\\
\emph{How close am I to the edge of the tube?} & Normalised errors: state divided by tube half-width & block 4\\
\emph{Which torque do I apply?} & A PD law plus a barrier term that grows near the edge, plus an adaptive term & blocks 5, 6\\
\bottomrule
\end{tabularx}
\end{center}

On the quadcopter four more questions appear: how to get the body rate when Gazebo only publishes the pose
(block 7), where the drone should be and therefore how it should be tilted (block 8), how a 60\,Hz digital loop
changes what the paper's continuous-time law may do (block 9), and how a torque is turned into four rotor speeds
(block 10).

\subsection{What is special about this controller}

Classical \emph{prescribed performance control} (PPC) keeps an error inside a tube by first transforming it into a
new variable that goes to infinity at the tube's edge, and then designing a controller for the transformed variable.
That works, but the controller then contains partial derivatives of the transformation and becomes long and coupled,
and it doubles in complexity when the angular velocity must be constrained as well. The paper's idea is simpler:
\begin{enumerate}[leftmargin=*]
  \item keep an ordinary PD controller ($-k_p\qv-k_d\vw$) that does the real work of turning the body;
  \item \emph{add} one auxiliary term whose gain is small in the middle of the tube and grows without bound near its
        edge, so the state is pushed back before it can leave;
  \item add one scalar adaptive estimate $\that$ that compensates the unknown disturbance and the torque that is lost
        to saturation;
  \item use a new \emph{finite-time} envelope that reaches its final width exactly at a chosen time $T_s$.
\end{enumerate}
The result is a law with a single line of algebra per axis, which is why the title calls it ``simple structure''.

\subsection{How every block is explained}

Each block of the two flow diagrams (\cref{sec:flows}) has its own subsection, always in the same order:

\begin{center}\small
\begin{tabularx}{\linewidth}{@{}l X@{}}
\toprule
\textbf{Part} & \textbf{What it gives you}\\
\midrule
Coloured equation box & The equations exactly as in the flow diagram, with the paper's equation numbers\\
Symbol table & Every symbol: meaning, size (dimension) and unit\\
Derivation & Where the equation comes from, step by step\\
\textcolor{cEnv}{\textbf{Toy example}} (orange box) & Small numbers you can check with a calculator\\
\textcolor{cOk}{\textbf{In our project}} (green box) & Where the equation lives in our code and what it does for our drone\\
Grey box & Which MATLAB step in \code{toy\_model/} prints the same numbers\\
\bottomrule
\end{tabularx}
\end{center}

\textbf{Three ways to read this guide.}
\begin{enumerate}[leftmargin=*]
  \item \emph{Ten-minute tour:} this section, the two flow diagrams and the recap table (\cref{sec:recap}).
  \item \emph{Understanding:} only the green ``In our project'' boxes and the toy examples, in order.
  \item \emph{Full learning:} everything in order, running the matching MATLAB step next to each block.
\end{enumerate}

% ---------------------------------------------------------------------
\subsection{Maths refresher}
\label{sec:refresher}
% ---------------------------------------------------------------------

Only the tools that are used later appear here, each with the smallest possible example.

\paragraph{Vectors, matrices, norms.} A vector $\bm a\in\R^3$ is a column of three numbers. Its length (Euclidean norm)
is $\|\bm a\|=\sqrt{a_1^2+a_2^2+a_3^2}$. The dot product is $\bm a^\top\bm b=\sum a_ib_i$.
\begin{toybox}
$\bm s=[0.13,\,-0.055,\,-0.135]^\top$: $\ \|\bm s\|=\sqrt{0.0169+0.003025+0.018225}=\sqrt{0.03815}=0.1953$. This exact
vector appears again in blocks 5 and 6.
\end{toybox}

\paragraph{The cross product as a matrix.} For $\bm z=[z_1,z_2,z_3]^\top$ the matrix
\begin{equation}\label{eq:skew}
\skw{\bm z}=\begin{bmatrix}0&-z_3&z_2\\ z_3&0&-z_1\\ -z_2&z_1&0\end{bmatrix}
\end{equation}
does the cross product: $\skw{\bm z}\bm y=\bm z\times\bm y$. It is skew-symmetric, $(\skw{\bm z})^\top=-\skw{\bm z}$, so
$\bm y^\top\skw{\bm z}\bm y=0$ for every $\bm y$ (a cross product is perpendicular to both factors). This single fact
makes the gyroscopic term disappear from the energy balance in Theorem~1.
\begin{toybox}
$\bm z=[1,0,0]^\top$ gives $\skw{\bm z}=\begin{bmatrix}0&0&0\\0&0&-1\\0&1&0\end{bmatrix}$. With $\bm y=[0,1,0]^\top$ the three rows
give $0$, $0$ and $1$, so $\skw{\bm z}\bm y=[0,0,1]^\top$, which is indeed $\bm e_x\times\bm e_y=\bm e_z$.
\end{toybox}

\paragraph{Quaternions.} A unit quaternion $\vq=[\qv^\top,q_4]^\top\in\R^4$, $\|\vq\|=1$, stores a rotation by an angle
$\phi$ about a unit axis $\bm n$:
\begin{equation}\label{eq:quat}
\qv=\bm n\sin\tfrac{\phi}{2}\in\R^3,\qquad q_4=\cos\tfrac{\phi}{2}\in\R .
\end{equation}
So $\vq=[0,0,0,1]^\top$ is ``no rotation'', and for small angles $\qv\approx\tfrac{\phi}{2}\bm n$: \emph{each component of
$\qv$ is about half the rotation angle about that axis}. That is why driving $\qv\to0$ means ``reach the target
orientation''. Quaternions have no singularities (unlike Euler angles at $\pm90^\circ$ pitch), but $\vq$ and $-\vq$ are
the same rotation, so every implementation keeps the sign continuous between samples. The product of two quaternions
(Hamilton product, written $\otimes$) composes two rotations, and the conjugate $\vq^*=[-\qv^\top,q_4]^\top$ is the
inverse rotation.
\begin{toybox}
A rotation of $60^\circ$ about $x$: $\qv=[\sin30^\circ,0,0]=[0.5,0,0]$, $q_4=\cos30^\circ=0.8660$; check
$0.5^2+0.866^2=1.000$. The paper's Part~1 initial attitude $\vq(0)=[0.2,-0.15,-0.25,0.9354]^\top$ has $q_4=0.9354$, so
$\phi=2\arccos(0.9354)=41.4^\circ$ away from the target.
\end{toybox}

\paragraph{Differential equations and RK4.} $\dot x=dx/dt$ is a rate of change. $\dot{\vx}=\bm f(\vx,\vu)$ tells how fast
the state changes, and a computer steps it forward with a small step $\Delta t$. Euler's method uses
$\vx_{k+1}=\vx_k+\Delta t\,\bm f(\vx_k)$. The fourth-order Runge--Kutta method (RK4) evaluates the slope four times
inside the step and averages them, $\vx_{k+1}=\vx_k+\tfrac{\Delta t}{6}(\bm k_1+2\bm k_2+2\bm k_3+\bm k_4)$; its error per
second of simulation shrinks like $\Delta t^4$. All our simulations use RK4 with $\Delta t=1$\,ms.

\paragraph{Saturation and sign functions.} $\sgn(x)$ is $+1$, $0$ or $-1$. The saturation function
$\sat(x)=\max(-1,\min(1,x))$ equals $x$ in $[-1,1]$ and is clipped outside. $\sat(x/\varsigma)$ is a smooth
replacement for $\sgn(x)$: it reaches $\pm1$ once $|x|\ge\varsigma$ and is linear in between (a ``boundary layer'').

\paragraph{Lyapunov functions.} To prove that a system settles, find an ``energy'' $V(\vx)\ge0$ that is zero only at the
target and show that it can only decrease along the motion, $\dot V\le0$. A ball in a bowl is the picture: height is
$V$, friction makes it go down, so the ball ends at the bottom.
\begin{toybox}
$\dot x=-2x$, $V=x^2$: $\dot V=2x\dot x=-4x^2\le0$, and $\dot V=0$ only at $x=0$. Indeed $x(t)=x(0)e^{-2t}\to0$.
\end{toybox}

\paragraph{$\mathcal L_2$ and Barbalat's lemma.} A signal is in $\mathcal L_2$ if $\int_0^\infty\|x\|^2dt<\infty$ (finite
``energy''). Barbalat's lemma: if a signal has finite energy \emph{and} its derivative is bounded (so it cannot spike
in ever thinner pulses), then the signal itself tends to zero. Theorem~1 of the paper uses exactly this argument.

\paragraph{Positive definite.} A symmetric matrix $M\succ0$ if $\vx^\top M\vx>0$ for every $\vx\ne0$. For a symmetric
$2\times2$ matrix this holds when both diagonal entries and the determinant are positive. $\lambda_{\min}(M)$ and
$\lambda_{\max}(M)$ are its smallest and largest eigenvalues, and $\lambda_{\min}\|\vx\|^2\le\vx^\top M\vx\le\lambda_{\max}\|\vx\|^2$.

\paragraph{Eigenvalues and discrete-time stability.} For $\vx_{k+1}=M\vx_k$ the state decays if every eigenvalue of $M$ has
$|\lambda|<1$ and blows up if one has $|\lambda|>1$. Block 9 uses this to decide how large a gain a 60\,Hz loop can take.
